\section{Traditional Interlocking Methods}

Traditional interlocking systems have historically ensured railway safety by preventing conflicting train movements through physical or hardwired logic. While effective in their time, their rigidity and dependence on manual or fixed infrastructure limit their suitability for modern, high-density networks.

\subsection{Mechanical Interlocking}

Mechanical interlocking was the earliest form, relying on lever frames connected by rods and wires to points and signals. Safety was guaranteed through physical locking between levers, preventing operators from setting incompatible routes. This design provided robust fail-safe protection but came with drawbacks: it was labor-intensive, required large installations of levers and linkages, and depended on manual operation, making it slow and impractical for expanding or complex networks.

\subsection{Relay-Based Interlocking}

Relay-based interlocking replaced mechanical rods with electrical circuits, encoding safety conditions through electromechanical relays. This allowed greater automation, reduced manual workload, and increased reliability compared to mechanical systems. However, the logic was hardwired, meaning any change to track layout or operating rules demanded costly rewiring. Large relay rooms for busy stations also consumed space and required constant maintenance, limiting scalability as railway operations grew more dynamic.